\begin{definition}[The energy-normalized vector Riesz transform]
For $f \in L^2(\gamma_n)$ set
\[
  R_\gamma f = D\, L^{-1/2}\Pi_0 f \in L^2(\gamma_n;\mathbb{R}^n),
\]
with $L^{-1/2}\Pi_0$ the pseudoinverse of \noderef{FracPow} and
$D = 2^{-1/2}\nabla$ the energy-normalized weak gradient of \noderef{GaussD}.
Its $j$th coordinate is $R_{\gamma,j} f = D_j L^{-1/2}\Pi_0 f$.  The arbitrary
$L^2$ input $f$ is never differentiated: one first forms the potential
$v = L^{-1/2}\Pi_0 f$, which lies in $\operatorname{Dom}(L^{1/2})$ and hence in
$W^{1,2}(\gamma_n)$ by \noderef{PseudoInverseSqrt} and
\noderef{SqrtDomainSobolev}, so its weak gradient exists in
$L^2(\gamma_n;\mathbb{R}^n)$.  Consequently $R_\gamma f$ is well defined for every
$f \in L^2(\gamma_n)$, with no differentiability assumption on $f$, and
$R_\gamma : L^2(\gamma_n) \to L^2(\gamma_n;\mathbb{R}^n)$ is linear.
\end{definition}
